\begin{definition}
For a \noderef{ScheideStructuralConfiguration} $C$, a \emph{Scheide structural
selection datum} consists of natural numbers $a,r$, a bijection
\[
 b:\operatorname{Fin}(r)\longrightarrow
   \{v:v\in\operatorname{ScheideExitVertices}(T,\varphi)\},
\]
and an injective retained-index map
$\rho:\operatorname{Fin}(r)\to\operatorname{Fin}(C.s)$.  The selected
color $C.\gamma(\rho(t))$ is missing at the exit underlying $b(t)$, and
among all selected colors it is the unique one missing there.  The datum
also contains a bijection
\[
 \upsilon:\operatorname{Fin}(a)\longrightarrow
 \{i\in\operatorname{Fin}(C.s):i\notin\operatorname{range}(\rho)\}.
\]
It records the equalities
$C.s=a+r$ and
$r=|\operatorname{ScheideExitVertices}(T,\varphi)|$, and records that each
old index belongs, uniquely and disjointly, either to the image of
$\upsilon$ or to the image of $\rho$.
\end{definition}
